% Checked inactive research fragment; no author input before integration.
\subsection{Homogeneous plurisubharmonic functions and complement indices}
\label{subsec:homogeneous-psh-complement-index}

Let $F\to B$ be a holomorphic bundle of rank $r\geq1$ on a connected compact complex manifold of dimension $n\geq1$. Write $F^\times$ for its complement of the zero section, and put $N=n+r$ and
\[
 U_F=\mathbb P(\mathcal O_B\oplus F)\setminus\mathbb P(\mathcal O_B).
\]
We use the smooth exhaustion definition of $q$-convexity in \cref{subsec:concave-neighbourhood-dual-norms}. Taking the reciprocal of any Hermitian squared norm gives an $N$-convex exhaustion. A homogeneous plurisubharmonic function forces this upper bound to be sharp.

\begin{lemma}\label{lem:universal-bundle-complement-index-bound}
The complement $U_F$ is $N$-convex.
\end{lemma}
\begin{proof}
Choose a smooth Hermitian metric $h$ and put $G(v)=h(v,v)$. By \cref{lem:reciprocal-bundle-norm-levi}, the vertical block of the Levi form of $G^{-1}$ has one positive eigenvalue at every point of $F^\times$. The block diagonal form therefore has at least one positive eigenvalue, irrespective of the horizontal block. The cutoff and extension at infinity in the proof of \cref{thm:bundle-curvature-complement-convexity} give a smooth exhaustion with this property outside a compact set. In dimension $N$, this is $N$-convexity.
\end{proof}

\begin{proposition}\label{prop:homogeneous-psh-sharp-complement-index}
Suppose that $F^\times$ admits a smooth positive plurisubharmonic function $G$ satisfying $G(cv)=|c|^2G(v)$ for every $c\in\mathbb C^*$ and $v\in F^\times$. Then the least smooth convexity index of $U_F$ is $N$.
\end{proposition}
\begin{proof}
It remains to exclude $(N-1)$-convexity, since every smaller convexity index would imply it. Suppose that a smooth exhaustion $\psi$ has at least two positive Levi eigenvalues outside a compact set $K\subset U_F$.

Fix a reference Hermitian norm on $F$. On its compact unit sphere bundle, $G$ has a positive minimum and a finite maximum. Homogeneity consequently compares $G$ with that squared norm uniformly on $F^\times$. Thus
\[
 \Sigma_t=\{G=t^2\},\qquad t>0,
\]
are compact smooth hypersurfaces, and multiplication by $t$ identifies $\Sigma_1$ with $\Sigma_t$. Smoothness follows also from the Euler identity: if $R_v=v$ is the real radial vector at $v$, then $\mathrm dG(R_v)=2G(v)>0$. Since $K$ stays away from the removed zero section, there is $\epsilon>0$ such that every $\Sigma_t$, $0<t<\epsilon$, is disjoint from $K$.

Put
\[
 M(t)=\max_{v\in\Sigma_1}\psi([1:tv]),\qquad 0<t<\epsilon.
\]
On any closed interval inside $(0,\epsilon)$, the radial derivatives of the smooth function under the maximum are uniformly bounded over the compact sphere. Hence $M$ is locally Lipschitz. Properness of the exhaustion implies $M(t)\to+\infty$ as $t\to0$: otherwise maximizing points along a sequence tending to zero would lie in one compact sublevel while accumulating on the removed section. Choose $0<t_0<t_1<\epsilon$ with $M(t_0)>M(t_1)$. The absolute continuity of a Lipschitz function and its integral derivative formula give a differentiability point $t\in(t_0,t_1)$ with $M'(t)<0$.

Choose $v\in\Sigma_1$ attaining $M(t)$, and set $p=tv$. For both signs of $h$, the comparison
\[
 M(t+h)\geq\psi([1:(t+h)v])
\]
is an equality at $h=0$. Differentiability of $M$ therefore gives $\frac{\mathrm d}{\mathrm dt}\psi([1:tv])=M'(t)$. Since $\psi$ has a constrained maximum on $\Sigma_t$ at $p$ and $\mathrm dG\ne0$, there is a real $\lambda$ with $\mathrm d\psi=\lambda\,\mathrm dG$ at $p$. The radial derivative gives
\[
 \lambda=\frac{M'(t)}{2t}<0.
\]
The function $\psi-\lambda G$ has zero differential at $p$ and nonpositive real Hessian on $T_p\Sigma_t$. Indeed, its restriction to the sphere is $\psi$ minus a constant, and its zero differential removes the second derivative of a curve realizing a tangent vector. On the complex tangent hyperplane $H=T_p\Sigma_t\cap\mathrm iT_p\Sigma_t$, this implies
\[
 \mathcal L_\psi|_H\leq\lambda\mathcal L_G|_H\leq0.
\]
The first inequality follows by adding the real Hessian inequalities in a tangent vector and its complex rotation; the second uses plurisubharmonicity and $\lambda<0$. A two-dimensional complex subspace on which $\mathcal L_\psi$ is positive definite would meet the complex codimension-one space $H$ nontrivially, a contradiction. Thus $\psi$ has at most one positive Levi eigenvalue at $p\notin K$. Combine this contradiction with \cref{lem:universal-bundle-complement-index-bound}.
\end{proof}

\begin{corollary}\label{cor:unitary-flat-exact-complement-index}
If $F$ admits a Hermitian metric with flat Chern connection, then the least smooth convexity index of $U_F$ is $N$.
\end{corollary}
\begin{proof}
In parallel holomorphic unitary frames, the squared norm is $G(v)=\sum_j|v_j|^2$, which is plurisubharmonic on the total space. Constant unitary transitions make it global, smooth, positive on $F^\times$ and homogeneous of degree two under complex scalars. Apply \cref{prop:homogeneous-psh-sharp-complement-index}.
\end{proof}

Fix $n\geq2$, a real $0<q<1$ and a complex $a\ne0$. On the scalar Hopf manifold $X_q=(\mathbb C^n\setminus\{0\})/\langle z\mapsto qz\rangle$, let $E_a$ be defined by $(z,v)\sim(qz,av)$, and put $b=-\log|a|/\log q$.

\begin{example}\label{ex:hopf-all-automorphy-complement-index}
The least smooth convexity index of $U_{E_a}$ is two when $0<|a|<1$, and $n+1$ when $|a|\geq1$. For every Gauduchon form $\gamma$ on $X_q$, the degree of $E_a$ is positive, zero or negative according as $|a|<1$, $|a|=1$ or $|a|>1$. Its de Rham first Chern class vanishes in all cases.
\end{example}
\begin{proof}
The squared norm $G(z,v)=|z|^{2b}|v|^2$ descends, since $q^{2b}|a|^2=1$. Its normalized curvature is $-b\alpha$, with $\alpha=\ddc\log|z|^2$ semipositive of rank $n-1$, as calculated in \cref{ex:hopf-exact-complement-index}. The integral of $\alpha\wedge\gamma^{n-1}$ is strictly positive, so the degree has the sign of $-b$, giving the three stated cases. The same scalar Hopf topology proves the vanishing de Rham class.

If $|a|<1$, then $b<0$. The explicit curvature again gives $2$-convexity by \cref{thm:bundle-curvature-complement-convexity}. Restriction to an embedded elliptic fibre of $X_q\to\mathbb P^{n-1}$ has flat Chern connection: the logarithmic metric coefficient there is $b\log|\zeta|^2$ plus a constant, hence pluriharmonic. A $1$-convex exhaustion would restrict to a proper strictly plurisubharmonic exhaustion outside a compact set of that line complement, contrary to \cref{lem:unitary-flat-complement-obstruction}. The phase of $a$ affects the gluing but not these metric and curvature calculations.

If $|a|\geq1$, then $b\geq0$. On the covering of $E_a^\times$, the function
\[
 \log G=b\log|z|^2+\log|v|^2
\]
is plurisubharmonic: the first summand has semipositive Levi form and the second is pluriharmonic for $v\ne0$. Consequently
\[
 \mathcal L_G=G\bigl(\mathcal L_{\log G}+\partial\log G\otimes\overline{\partial\log G}\bigr)\geq0.
\]
The descended $G$ is smooth, positive and homogeneous of degree two under complex scalars. Apply \cref{prop:homogeneous-psh-sharp-complement-index} with rank one and total dimension $n+1$.
\end{proof}
